\documentclass[12pt]{article}
\usepackage[utf8]{inputenc}
\usepackage[T1]{fontenc}
\usepackage[margin=1in]{geometry}
\usepackage{setspace}
\usepackage{xcolor}
\usepackage[colorlinks=true, urlcolor=blue!60!black, linkcolor=blue!60!black]{hyperref}
\setlength{\parindent}{2em}
\setlength{\parskip}{1.25em}
\renewcommand{\baselinestretch}{1.0}
\begin{document}
\begin{center}
{\LARGE \textit{National Federation of Independent Business v. Sebelius}, 567 U.S. 519 (2012)\par}
\vspace{0.8em}
{\large Case Brief\par}
\end{center}
\vspace{1em}

\section{Memory Jogger}
In a fractured landmark decision, Chief Justice Roberts joined the four liberal Justices to uphold the Affordable Care Act's individual mandate as a valid exercise of Congress's taxing power---even while five Justices agreed it exceeded the Commerce Clause---and joined six other Justices to hold that threatening States with the loss of all existing Medicaid funding to enforce the Act's Medicaid expansion was unconstitutionally coercive under the Spending Clause.

\section{Detailed Case Facts}
In March 2010, Congress enacted the Patient Protection and Affordable Care Act (ACA), Pub. L. No. 111--148, 124 Stat. 119, a 900-page, 10-title statute designed to expand health insurance coverage and reduce health care costs (\href{https://www.courtlistener.com/opinion/809122/national-federation-of-independent-business-v-sebelius/}{CourtListener}; \href{https://supreme.justia.com/cases/federal/us/567/519/}{Justia}). At the time of enactment, Americans spent \$2.5 trillion annually on health care (17.6\% of the national economy), roughly 50 million U.S. residents lacked health insurance, and federal and state emergency-care laws (such as EMTALA, 42 U.S.C. \S~1395dd) required hospitals to treat uninsured patients regardless of ability to pay, generating \$43 billion in uncompensated care in 2008 that Congress found raised insured families' annual premiums by over \$1,000 on average (42 U.S.C. \S~18091(2)(F)) (\href{https://supreme.justia.com/cases/federal/us/567/519/}{Justia}).

To expand access for people with preexisting conditions, the ACA enacted ``guaranteed-issue'' and ``community-rating'' reforms (42 U.S.C. \S\S~300gg, 300gg--1, 300gg--3, 300gg--4), prohibiting insurers from denying coverage or charging higher premiums based on health status. Because those reforms alone create an adverse-selection risk---encouraging healthy individuals to wait until they fall ill to buy coverage---Congress paired them with two challenged provisions (\href{https://supreme.justia.com/cases/federal/us/567/519/}{Justia}):

\begin{enumerate}
\item \textbf{The Individual Mandate (26 U.S.C. \S~5000A):} Beginning in 2014, most non-exempt Americans were required to maintain ``minimum essential'' health insurance coverage or pay a ``shared responsibility payment''---labeled a ``penalty'' in the statute---to the Internal Revenue Service (IRS) with their annual federal income tax returns (\S~5000A(b)(1), (g)(1)). By 2016, the exaction was set at 2.5\% of household income, with a floor of \$695 and a ceiling capped at the national average annual premium for a qualifying bronze-level plan (\S~5000A(c)). Individuals below the federal income tax filing threshold were exempt from the payment (\S~5000A(e)(2)), and the IRS was expressly barred from enforcing it through criminal prosecutions, liens, or levies (\S~5000A(g)(2)) (\href{https://supreme.justia.com/cases/federal/us/567/519/}{Justia}).

\item \textbf{The Medicaid Expansion (42 U.S.C. \S\S~1396a(a)(10)(A)(i)(VIII), 1396c, 1396d(y)(1)):} Enacted in 1965, traditional Medicaid provided matching federal funds (50\% to 83\% of program costs) to States to cover discrete categories of needy residents: pregnant women, children, needy families, the blind, the elderly, and the disabled. By 1982, every State participated, and federal Medicaid grants accounted for over 10\% of most States' total budgets (\href{https://supreme.justia.com/cases/federal/us/567/519/}{Justia}). The ACA required participating States to expand Medicaid by 2014 to cover virtually all nonelderly adults with incomes up to 133\% of the federal poverty level, with the federal government funding 100\% of newly eligible costs through 2016 and phasing down to 90\% thereafter (\S~1396d(y)(1)). Under existing 42 U.S.C. \S~1396c, however, if a State declined to adopt the expansion, the Secretary of Health and Human Services could withhold \textit{all} of that State's federal Medicaid funds, including pre-existing funding (\href{https://supreme.justia.com/cases/federal/us/567/519/}{Justia}).
\end{enumerate}

\section{Procedural History}
On the day President Barack Obama signed the ACA into law, Florida and twelve other States sued Secretary of Health and Human Services Kathleen Sebelius and federal agencies in the U.S. District Court for the Northern District of Florida; thirteen additional States (totaling twenty-six States), individual plaintiffs, and the National Federation of Independent Business (NFIB) later joined the suit (\href{https://supreme.justia.com/cases/federal/us/567/519/}{Justia}; \href{https://www.law.cornell.edu/supremecourt/text/11-393}{Cornell LII}). On January 31, 2011, District Judge Roger Vinson rejected the challenge to the Medicaid expansion but held that the individual mandate exceeded Congress's Article I powers, concluded the mandate was inseverable from the rest of the statute, and invalidated the entire ACA, 780 F. Supp. 2d 1256 (N.D. Fla. 2011) (\href{https://supreme.justia.com/cases/federal/us/567/519/}{Justia}).

On appeal, the U.S. Court of Appeals for the Eleventh Circuit affirmed in part and reversed in part, 648 F.3d 1235 (11th Cir. 2011). The panel unanimously upheld the Medicaid expansion under the Spending Clause and unanimously rejected the federal government's taxing-power defense of the individual mandate; a 2--1 majority (with Judge Marcus dissenting) held that the mandate exceeded the Commerce Clause, but reversed the District Court on severability, striking only the mandate while leaving the remainder of the ACA intact (\href{https://supreme.justia.com/cases/federal/us/567/519/}{Justia}). Meanwhile, a circuit split developed: the Sixth Circuit (\textit{Thomas More Law Center v. Obama}, 651 F.3d 529 (6th Cir. 2011)) and D.C. Circuit (\textit{Seven-Sky v. Holder}, 661 F.3d 1 (D.C. Cir. 2011)) upheld the mandate under the Commerce Clause, while the Fourth Circuit (\textit{Liberty Univ., Inc. v. Geithner}, 671 F.3d 391 (4th Cir. 2011)) held that the Tax Anti-Injunction Act, 26 U.S.C. \S~7421(a), barred pre-enforcement review (\href{https://supreme.justia.com/cases/federal/us/567/519/}{Justia}).

The Supreme Court granted certiorari on November 14, 2011, consolidating Nos. 11--393 (\textit{NFIB v. Sebelius}), 11--398 (\textit{HHS v. Florida}), and 11--400 (\textit{Florida v. HHS}), appointed court-invited \textit{amici curiae} to argue the Anti-Injunction Act and severability issues, heard an extraordinary three days of oral argument (March 26--28, 2012), and issued its decision on June 28, 2012 (\href{https://supreme.justia.com/cases/federal/us/567/519/}{Justia}; \href{https://www.law.cornell.edu/supremecourt/text/11-393}{Cornell LII}).

\section{Judicial Votes}
The Court's judgment---affirming in part and reversing in part the Eleventh Circuit---rested on shifting coalitions across four distinct constitutional and statutory issues (\href{https://www.law.cornell.edu/supremecourt/text/11-393}{Cornell LII}; \href{https://wustllawreview.org/wp-content/uploads/2021/10/91.1.1.pdf}{Solum, Wash. U. L. Rev.}):

\begin{itemize}
\item \textbf{Anti-Injunction Act (9--0):} All nine Justices agreed that 26 U.S.C. \S~7421(a) did not bar pre-enforcement review. Chief Justice Roberts delivered the Opinion of the Court in Part II, joined by Justices Ginsburg, Breyer, Sotomayor, and Kagan; Justices Scalia, Kennedy, Thomas, and Alito agreed in their joint dissent.
\item \textbf{Individual Mandate---Commerce \& Necessary and Proper Clauses (5--4 against federal power, no single opinion):} Chief Justice Roberts (Parts III--A, III--B, and III--D, writing for himself alone) and the four joint dissenters (Scalia, Kennedy, Thomas, and Alito, JJ.) concluded that Congress lacked authority under Art. I, \S~8, cls. 3 and 18 to compel uninsured individuals to buy health insurance. Justice Thomas also filed a brief solo dissent rejecting the ``substantial effects'' test altogether. Justices Ginsburg, Breyer, Sotomayor, and Kagan (Ginsburg, J., Parts I--IV) would have upheld the mandate under both Clauses.
\item \textbf{Individual Mandate---Taxing Power (5--4 upholding the mandate):} Chief Justice Roberts delivered the Opinion of the Court in Parts I, II, and III--C, joined by Justices Ginsburg, Breyer, Sotomayor, and Kagan, upholding 26 U.S.C. \S~5000A under Art. I, \S~8, cl. 1. Justices Scalia, Kennedy, Thomas, and Alito dissented.
\item \textbf{Medicaid Expansion---Spending Clause Coercion (7--2):} Seven Justices concluded that threatening to strip all existing Medicaid funds under 42 U.S.C. \S~1396c was unconstitutionally coercive: Chief Justice Roberts (Part IV--A), joined by Justices Breyer and Kagan (3-Justice plurality), plus the four joint dissenters (Scalia, Kennedy, Thomas, and Alito, JJ.). Justices Ginsburg and Sotomayor (Ginsburg, J., Part V) dissented on coercion.
\item \textbf{Medicaid Expansion---Remedy and Severability (5--4):} Chief Justice Roberts (Part IV--B), joined by Justices Breyer and Kagan, limited the remedy to barring the Secretary from applying \S~1396c to withdraw pre-existing Medicaid funds, leaving the expansion optional and the rest of the ACA intact. Given the 7-Justice coercion ruling, Justices Ginsburg and Sotomayor agreed with Part IV--B's narrow severability remedy (5 votes total for the remedy). Justices Scalia, Kennedy, Thomas, and Alito would have struck down the Medicaid expansion and the entire ACA as inseverable.
\end{itemize}

\section{Holding}
\subsection{Judgment}
The judgment of the Eleventh Circuit was \textbf{affirmed in part and reversed in part}: the Court reversed the Eleventh Circuit's invalidation of the individual mandate (upholding it as a tax) and reversed the Eleventh Circuit's ruling on the Medicaid expansion (holding the threatened withdrawal of existing Medicaid funding unconstitutionally coercive, while severing that enforcement mechanism so the expansion became optional for States and the remainder of the ACA survived) (\href{https://supreme.justia.com/cases/federal/us/567/519/}{Justia}).

\subsection{Rule Announced}
\begin{enumerate}
\item \textbf{Anti-Injunction Act:} Because Congress labeled the ACA's shared responsibility payment a ``penalty'' rather than a ``tax,'' the statutory bar of 26 U.S.C. \S~7421(a) does not preclude pre-enforcement judicial review.
\item \textbf{Taxing Power (Art. I, \S~8, cl. 1):} Under a functional constitutional inquiry---disregarding statutory labels---an exaction that imposes a non-prohibitory financial burden, contains no scienter requirement, is collected solely by the IRS through normal tax administration without criminal sanctions, raises revenue, and leaves citizens a lawful choice to pay in lieu of purchasing insurance is a valid exercise of Congress's taxing power and is not a ``direct'' tax requiring apportionment under Art. I, \S~9, cl. 4.
\item \textbf{Spending Clause Coercion (Art. I, \S~8, cl. 1):} Congress crosses the line from permissible persuasion to unconstitutional coercion when it threatens to terminate an independent, pre-existing federal grant accounting for a massive share of state budgets (over 10\%) to compel States to accept a new program that effects a ``shift in kind, not merely degree.''
\end{enumerate}

\subsection{Holding vs. Dicta}
The status of the Commerce Clause analysis generated an enduring doctrinal debate (\href{https://wustllawreview.org/wp-content/uploads/2021/10/91.1.1.pdf}{Solum, Wash. U. L. Rev.}):
\begin{itemize}
\item \textbf{Binding Holdings:} The Anti-Injunction Act ruling (Part II, 5-Justice majority), the Taxing Power ruling (Part III--C, 5-Justice majority), and the Spending Clause coercion holding and narrow severability remedy (Part IV, 3-Justice plurality controlling under \textit{Marks v. United States}, 430 U.S. 188 (1977), reinforced by 7 votes on coercion and 5 votes on remedy) are indisputably binding holdings.
\item \textbf{Contested Status of the Commerce Clause Analysis (Parts III--A, III--B, III--D):} Justice Ginsburg argued that because the mandate was valid as a tax, Chief Justice Roberts's solo rejection of the Commerce Clause theory was unnecessary dicta. Chief Justice Roberts explicitly countered in Part III--D that deciding the Commerce Clause issue was a \textit{necessary logical predicate} to his judgment: because the most natural reading of \S~5000A was a command to buy insurance, he could adopt the ``saving construction'' treating it as a tax only after first holding that the command reading was unconstitutional under the Commerce Clause (\href{https://supreme.justia.com/cases/federal/us/567/519/}{Justia}). Scholars remain divided on whether Roberts's Commerce Clause analysis combined with the four-Justice joint dissent creates binding horizontal or vertical precedent against future activity-compelling mandates (\href{https://wustllawreview.org/wp-content/uploads/2021/10/91.1.1.pdf}{Solum, Wash. U. L. Rev.}).
\end{itemize}

\section{Analysis of Opinions}
\subsection{Opinion of the Court and Controlling Plurality (Roberts, C.J.)}
Chief Justice Roberts structured his opinion around the foundational premise that the Federal Government is one of limited, enumerated powers, whereas the general police power remains with the States (\href{https://supreme.justia.com/cases/federal/us/567/519/}{Justia}).

In \textbf{Part II} (for the Court), Roberts held that the Anti-Injunction Act did not apply because both the AIA and the ACA are creatures of Congress, making Congress's choice of the word ``penalty'' rather than ``tax'' decisive for statutory jurisdiction.

In \textbf{Part III--A} (for himself alone), Roberts rejected the Commerce Clause and Necessary and Proper Clause justifications. Textually and historically, the power to ``regulate Commerce'' presupposes existing commercial \textit{activity}; compelling inactive individuals to enter a market to purchase an unwanted product---like ordering citizens to buy broccoli to lower health care costs---would grant Congress a boundless police power to regulate citizens ``from cradle to grave.'' Even under \textit{Wickard v. Filburn}, 317 U.S. 111 (1942), the farmer was actively growing wheat. Nor was the mandate ``proper'' under the Necessary and Proper Clause, because rather than carrying an enumerated power into execution, it created the very commercial predicate needed to exercise that power (\href{https://supreme.justia.com/cases/federal/us/567/519/}{Justia}).

In \textbf{Part III--B} (solo) and \textbf{Part III--C} (for a 5-Justice majority), Roberts invoked constitutional avoidance: if a statute is ``fairly possible'' to read in a constitutional manner, courts must adopt that saving construction (\textit{Crowell v. Benson}, 285 U.S. 22, 62 (1932)). Applying the functional test of \textit{Bailey v. Drexel Furniture Co.}, 259 U.S. 20 (1922), Roberts concluded that \S~5000A operated as a valid tax: the payment was modest (capped below the cost of insurance), lacked a scienter requirement, was collected solely by the IRS without criminal penalties, raised roughly \$4 billion annually, and imposed no legal stigma or unlawfulness on citizens who chose to pay rather than buy insurance. It was not a direct tax (neither a capitation nor a property tax) and thus did not require apportionment. Crucially, unlike the commerce power, the taxing power does not allow Congress to criminalize inactivity; it only requires paying money to the Treasury (\href{https://supreme.justia.com/cases/federal/us/567/519/}{Justia}).

In \textbf{Part IV} (joined by Breyer and Kagan, JJ.), Roberts applied the contract analogy of Spending Clause jurisprudence (\textit{Pennhurst State School and Hospital v. Halderman}, 451 U.S. 1 (1981); \textit{South Dakota v. Dole}, 483 U.S. 203 (1987)). Whereas \textit{Dole} threatened less than 0.5\% of South Dakota's budget (5\% of highway funds), \S~1396c threatened over 10\% of a State's entire budget---not ``relatively mild encouragement,'' but ``a gun to the head'' and ``economic dragooning.'' Because the ACA transformed Medicaid from a categorical program for discrete needy groups into a universal low-income health plan (a ``shift in kind, not merely degree''), Congress's reservation of the right to ``alter'' or ``amend'' Medicaid (42 U.S.C. \S~1304) gave States no fair notice. In \textbf{Part IV--B}, relying on the Medicaid Act's severability clause (42 U.S.C. \S~1303), Roberts remedied the violation narrowly by barring the Secretary from withdrawing pre-existing Medicaid funds from non-participating States, leaving the rest of the ACA intact (\href{https://supreme.justia.com/cases/federal/us/567/519/}{Justia}).

\subsection{Concurrence in Part, Concurrence in the Judgment in Part, and Dissent in Part (Ginsburg, J.)}
Justice Ginsburg---joined by Breyer, Sotomayor, and Kagan, JJ., in Parts I--IV, and by Sotomayor, J., in Part V---joined Roberts's Anti-Injunction Act and Taxing Power holdings but sharply criticized his Commerce Clause and Spending Clause reasoning (\href{https://supreme.justia.com/cases/federal/us/567/519/}{Justia}; \href{https://en.wikisource.org/wiki/National_Federation_of_Independent_Business_v._Sebelius/Concurrence-dissent_Ginsburg}{Wikisource}).

On the \textbf{Commerce Clause (Parts I--III)}, Ginsburg called Roberts's activity/inactivity distinction a ``crabbed reading'' that ``harks back to the era'' of pre-1937 formalism. Unlike cars or broccoli, health care is a market virtually every person inevitably and unpredictably enters, where emergency care is legally guaranteed regardless of ability to pay, generating massive cost-shifting and interstate collective-action problems. Choosing to self-insure is an economic financing decision within an existing market. Under \textit{Gonzales v. Raich}, 545 U.S. 1 (2005), the mandate was also plainly ``necessary and proper'' as an essential component of the concededly valid guaranteed-issue and community-rating scheme, without which insurance markets suffer an adverse-selection ``death spiral.''

On the \textbf{Spending Clause (Part V, joined by Sotomayor, J.)}, Ginsburg argued that the ACA merely expanded a single, existing cooperative program that Congress had amended dozens of times under \S~1304, with the federal government paying 90\% to 100\% of new costs. She warned that the plurality's coercion test lacked judicially manageable standards. Given the 7-Justice coercion holding, however, she joined Roberts's Part IV--B remedy to conserve the statute rather than destroy it (\href{https://supreme.justia.com/cases/federal/us/567/519/}{Justia}).

\subsection{Joint Dissent (Scalia, Kennedy, Thomas, and Alito, JJ.)}
In an unusual unsigned joint dissent, Justices Scalia, Kennedy, Thomas, and Alito agreed with Roberts that the individual mandate exceeded the Commerce and Necessary and Proper Clauses and that the Medicaid expansion was unconstitutionally coercive, but would have struck down the entire ACA (\href{https://en.wikisource.org/wiki/National_Federation_of_Independent_Business_v._Sebelius/Dissent_Scalia,_Kennedy,_Thomas,_and_Alito}{Wikisource}).

On the \textbf{Taxing Power}, the joint dissent accused the majority of ``verbal wizardry'' and judicial legislation: Congress repeatedly called \S~5000A a ``requirement'' and a ``penalty''---not a tax---and imposed a legal duty to buy insurance. Rewriting a regulatory penalty into a tax violated the structural accountability of taxation. On \textbf{Severability}, the dissenters argued that the individual mandate and Medicaid expansion were the two structural pillars of the ACA's health-insurance architecture; without them, Congress would not have enacted the remaining provisions, requiring invalidation of the entire Act (\href{https://en.wikisource.org/wiki/National_Federation_of_Independent_Business_v._Sebelius/Dissent_Scalia,_Kennedy,_Thomas,_and_Alito}{Wikisource}).

\subsection{Separate Dissent (Thomas, J.)}
Justice Thomas added a one-paragraph solo dissent reiterating his longstanding originalist objection that the Court's ``substantial effects'' test under the Commerce Clause is ``inconsistent with the original understanding of Congress' powers and with this Court's early Commerce Clause cases'' (\href{https://supreme.justia.com/cases/federal/us/567/519/}{Justia}).

\section{Examples: Future Applications}
\begin{enumerate}
\item \textbf{Same-Side (Valid under Taxing Power): Federal Climate-Resilience Insurance Surcharge.} Congress enacts a statute requiring homeowners in federally designated flood zones to carry private flood insurance or pay a modest annual \$400 ``climate resilience assessment'' collected by the IRS on Form 1040, with no criminal sanctions, no scienter element, and an exemption for low-income filers. \textit{Why Same-Side:} Even if compelling uninsured homeowners to enter the insurance market exceeds the Commerce Clause under Roberts's activity/inactivity line, the exaction satisfies every functional criterion of \textit{NFIB} Part III--C: it raises revenue, is non-prohibitory, is administered solely by the IRS, and leaves homeowners a lawful choice to pay the tax instead of buying insurance.

\item \textbf{Same-Side (Invalid under Spending Clause): Threatening All Federal Education Aid to Mandate a New Universal Pre-K Program.} Congress creates a brand-new universal preschool program for all three- and four-year-olds and directs the Secretary of Education to withhold 100\% of a State's existing Title I K--12 federal education funding---representing roughly 10\% of state education budgets---if the State refuses to adopt the new program. \textit{Why Same-Side:} Like the ACA's Medicaid expansion, this couples a ``shift in kind'' (a new universal program distinct from existing K--12 grants) with the threatened termination of massive, entrenched pre-existing federal funds, crossing from persuasion into ``economic dragooning'' under Part IV--A.

\item \textbf{Opposite-Side (Invalid---Cannot Be Saved as a Tax): Criminal Mandate to Purchase Electric Vehicles.} Congress passes a law ordering every licensed U.S. driver to purchase an electric vehicle by 2030, enforced by the Department of Transportation through a \$15,000 fine and up to six months' imprisonment for willful noncompliance. \textit{Why Opposite-Side:} First, under the five-Justice Commerce Clause reasoning in \textit{NFIB}, compelling inactive individuals to purchase a product they do not want exceeds the commerce power. Second, unlike \S~5000A, this mandate cannot be saved as a tax under \textit{Drexel Furniture} and \textit{NFIB} Part III--C because it imposes a prohibitory fine, requires willful scienter, is enforced by a regulatory agency rather than the IRS, and attaches criminal liability.

\item \textbf{Opposite-Side (Valid under Spending Clause): Conditioning 4\% of Federal Highway Funds on Automated Speed-Enforcement Cameras.} Congress amends the Federal-Aid Highway Act to withhold 4\% of a State's annual federal highway construction grants (less than 0.3\% of the average State's total budget) unless the State installs automated speed-enforcement cameras in highway work zones. \textit{Why Opposite-Side:} Unlike the Medicaid expansion's ``gun to the head,'' this mirrors \textit{South Dakota v. Dole}: the condition is germane to highway safety, involves a small fraction of a single grant rather than a threat to over 10\% of a State's overall budget, and leaves States a genuine choice in theory and in fact.

\item \textbf{Fence-Sitter: Zero-Dollar Mandate or Hybrid Tax-Penalty on Unvaccinated Adults.} Suppose Congress enacts a ``public health shared responsibility'' statute requiring adults to receive an annual respiratory vaccine or pay a \$2,500 IRS-collected assessment (roughly 50 times the \$50 cost of the vaccine), with no criminal penalty and no scienter requirement, \textit{or} later reduces the assessment to \$0 (as Congress did to the ACA mandate in the 2017 Tax Cuts and Jobs Act, litigated in \textit{California v. Texas}, 593 U.S. 659 (2021)). \textit{Why Fence-Sitter:} A \$2,500 exaction for foregoing a \$50 vaccine sits squarely on the unresolved boundary \textit{NFIB} declined to fix---where a levy collected by the IRS without scienter nonetheless becomes so steep and ``prohibitory'' under \textit{Drexel Furniture} that it loses its character as a tax and becomes an unconstitutional penalty. Conversely, a \$0 exaction raises zero revenue, stripping away ``the essential feature of any tax'' under \textit{NFIB} Part III--C while simultaneously raising Article III standing barriers against challengers.
\end{enumerate}

\section{Critique}
\subsection{Originalist and Textualist Perspectives}
Originalist and textualist scholars reacted with sharply divided assessments. On the Commerce Clause, scholars such as Randy Barnett celebrated \textit{NFIB} as a structural victory that halted the post-1937 trajectory of unlimited federal power and destabilized the ``New Deal Settlement,'' shifting the ``constitutional gestalt'' toward enforceable enumerated limits (\href{https://wustllawreview.org/wp-content/uploads/2021/10/91.1.1.pdf}{Solum, Wash. U. L. Rev.}; \href{https://www.cato.org/outside-articles/eastwood-talking-chair-good-bad-ugly-obamacare-ruling}{Cato Institute}). Conversely, textualists---including then-Professor Amy Coney Barrett in a 2017 book review---sharply criticized Chief Justice Roberts's taxing-power ``saving construction'' as pushing the statutory text ``beyond its plausible meaning'' to rescue the ACA, echoing the joint dissent's charge that calling a statutory ``penalty'' a ``tax'' undermined textual fidelity and democratic accountability (\href{https://reason.com/volokh/2020/09/26/what-did-professor-amy-coney-barrett-say-about-nfib-v-sebelius/}{Volokh Conspiracy}; \href{https://texaslawreview.org/facial-challenges-saving-constructions-and-statutory-severability/}{Texas Law Review}). On the Spending Clause, conservative scholar Michael S. Greve praised the Court for finally giving teeth to anti-coercion limits, yet cautioned that Roberts's fact-bound ``shift in kind'' and budget-percentage formula left cooperative federalism's broader structural pathologies unresolved (\href{https://journals.law.harvard.edu/jlpp/wp-content/uploads/sites/90/2014/01/37_1_83_Greve.pdf}{Greve, Harv. J.L. \& Pub. Pol'y}).

\subsection{Progressive and Doctrinal Perspectives}
Progressive and doctrinal scholars leveled three major critiques. First, commentators argued that the activity/inactivity distinction is conceptually arbitrary and mirrors the discredited formalisms of the \textit{Lochner} era and the state-action doctrine, because economic choices (such as self-insuring) can always be redescribed as either action or inaction (\href{https://harvardlawreview.org/wp-content/uploads/2014/02/vol127_nfib_v_sebelius.pdf}{Harvard Law Review Note, 127 Harv. L. Rev. 1174}). Second, scholars such as Andrew Coan and David S. Schwartz have challenged the originalist pedigree of strict ``enumerationism'' itself, arguing that founding-era public meaning regarding the General Welfare, Necessary and Proper, and Commerce Clauses was contested and indeterminate (\href{https://ilr.law.uiowa.edu/sites/ilr.law.uiowa.edu/files/2024-03/ILR-109-Coan-Schwartz_0.pdf}{Coan \& Schwartz, Iowa L. Rev.}). Third, critics of ``active avoidance'' contended that Roberts's methodology allowed the Court to announce major new constitutional restrictions on the Commerce Clause without bearing the political cost of striking down the statute, while the Medicaid coercion holding created a ``coverage gap'' affecting millions of low-income residents in non-expanding States and ushered in an era of state-strengthening Roberts Court federalism (\href{https://harvardlawreview.org/print/vol-138/federalism-rebalancing-and-the-roberts-court-a-departure-from-historical-patterns/}{Harvard Law Review, 138 Harv. L. Rev.}).

\subsection{Underlying Competing Values}
\textit{NFIB} exposes the core tension of American constitutional federalism: preserving \textbf{limited, enumerated federal authority and state sovereign autonomy} (which diffuse power, protect individual liberty, and maintain clear lines of political accountability) against enabling \textbf{effective national solutions to interstate collective-action problems} (such as universal health-care financing, which individual States struggle to solve alone without triggering adverse migration and economic disadvantage).

\section{Key Quotations}
\begin{quote}
``The Framers gave Congress the power to regulate commerce, not to compel it, and for over 200 years both our decisions and Congress's actions have reflected this understanding. There is no reason to depart from that understanding now.''
\end{quote}
---Chief Justice Roberts, Part III--A--1, 567 U.S. at 555 (\href{https://supreme.justia.com/cases/federal/us/567/519/}{Justia}; \href{https://www.courtlistener.com/opinion/809122/national-federation-of-independent-business-v-sebelius/}{CourtListener}).

\begin{quote}
``The Affordable Care Act's requirement that certain individuals pay a financial penalty for not obtaining health insurance may reasonably be characterized as a tax. Because the Constitution permits such a tax, it is not our role to forbid it, or to pass upon its wisdom or fairness.''
\end{quote}
---Chief Justice Roberts, Opinion of the Court, Part III--C, 567 U.S. at 574 (\href{https://supreme.justia.com/cases/federal/us/567/519/}{Justia}; \href{https://www.courtlistener.com/opinion/809122/national-federation-of-independent-business-v-sebelius/}{CourtListener}).

\begin{quote}
``In this case, the financial `inducement' Congress has chosen is much more than `relatively mild encouragement'---it is a gun to the head. \ldots\ The threatened loss of over 10 percent of a State's overall budget, in contrast, is economic dragooning that leaves the States with no real option but to acquiesce in the Medicaid expansion.''
\end{quote}
---Chief Justice Roberts, Plurality Opinion, Part IV--A, 567 U.S. at 581--82 (\href{https://supreme.justia.com/cases/federal/us/567/519/}{Justia}; \href{https://www.courtlistener.com/opinion/809122/national-federation-of-independent-business-v-sebelius/}{CourtListener}).

\begin{quote}
``The Chief Justice's crabbed reading of the Commerce Clause harks back to the era in which the Court routinely thwarted Congress' efforts to regulate the national economy in the interest of those who labor to sustain it.''
\end{quote}
---Justice Ginsburg, Concurring in Part, Concurring in the Judgment in Part, and Dissenting in Part, Part I, 567 U.S. at 589 (\href{https://supreme.justia.com/cases/federal/us/567/519/}{Justia}; \href{https://www.courtlistener.com/opinion/809122/national-federation-of-independent-business-v-sebelius/}{CourtListener}).

\begin{quote}
``The Act before us here exceeds federal power both in mandating the purchase of health insurance and in denying non-consenting States all Medicaid funding. These parts of the Act are central to its design and operation, and all the Act's other provisions would not have been enacted without them. In our view it must follow that the entire statute is inoperative.''
\end{quote}
---Justices Scalia, Kennedy, Thomas, and Alito, Joint Dissent, 567 U.S. at 648--49 (\href{https://en.wikisource.org/wiki/National_Federation_of_Independent_Business_v._Sebelius/Dissent_Scalia,_Kennedy,_Thomas,_and_Alito}{Wikisource}; \href{https://www.courtlistener.com/opinion/809122/national-federation-of-independent-business-v-sebelius/}{CourtListener}).

\end{document}
